\documentclass[11pt]{article}
\usepackage[a4paper,margin=1in]{geometry}
\usepackage{amsmath,amssymb,amsthm}
\usepackage{booktabs}
\usepackage{graphicx}
\usepackage{xcolor}
\usepackage{hyperref}
\hypersetup{colorlinks=true,linkcolor=blue,urlcolor=blue,citecolor=blue}
\newcommand{\rt}[1]{\texttt{#1}}
\newcommand{\fn}[1]{\texttt{#1}}
\providecommand{\ket}[1]{|#1\rangle}
\providecommand{\bra}[1]{\langle #1|}
\newcommand{\Fix}{\operatorname{Fix}}
\newcommand{\cmark}{\ensuremath{\checkmark}}
\newcommand{\xmark}{\ensuremath{\times}}
\newcommand{\ifgraphic}[2]{\IfFileExists{#1}{\includegraphics[width=#2]{#1}}%
  {\fbox{\parbox{#2}{\centering\small figure \texttt{\detokenize{#1}} pending
  (regenerate after the sweep)}}}}

\theoremstyle{plain}
\newtheorem{theorem}{Theorem}
\newtheorem{proposition}{Proposition}
\newtheorem{corollary}{Corollary}
\theoremstyle{definition}
\newtheorem{definition}{Definition}

\title{\textbf{Report R8 --- C150 at open boundaries: an exactly solved unitary
QCA, and exactly what it leaves open}\\
\large Krylov sector analysis of quantum cellular automata}
\author{QCA Sector Analysis project (Tier 1)}
\date{2026-07-26 \quad\small(engine\_version \texttt{1a.1})}

\begin{document}
\maketitle

\begin{abstract}
C150 --- the quantum elementary cellular automaton with Wolfram number $150$,
local symbol table $(\rt{I},\rt{V},\rt{V},\rt{I})$, written \rt{W150} elsewhere
in this series --- is the member of the $16$-rule unitary family whose sector
structure we can solve in closed form at open boundaries. We do so here, and
then measure precisely how much of the dynamics that solution does \emph{not}
determine. Three results. (i) \emph{Reduction.} For any unitary rule the sector
partition equals the connected components of the single-flip graph
$x \sim x\oplus 2^i$ over the sites at which the Hadamard fires; the brick-wall
ordering drops out. Verified against the streamed engine on all $16$ unitary
rules, both boundary conventions, $N=3\ldots12$ ($320$ units, no mismatch).
(ii) \emph{Solution.} In the bond variables $b_j = x_j \oplus x_{j+1}$ the map is
a bijection onto the even-weight strings of length $N+1$ and the elementary move
is a hard-core \emph{hop} of one domain wall. Hence the conserved charge is the
open Ising domain-wall number $Q=\sum_j (1-Z_jZ_{j+1})/2$, the sectors are its
eigenspaces, $|S_w| = \binom{N+1}{w}$ for even $w$, the sector count is
$\lfloor (N{+}1)/2\rfloor + 1$ and $D_{\max} = \max_{w\,\mathrm{even}}
\binom{N+1}{w}$. Because the sector count is \emph{linear} in $N$, C150 is
symmetry-resolved and not fragmented: it is the natural null model for the
family. The exponential base of $D_{\max}$ is exactly $2$ with an $N^{-1/2}$
prefactor, which supersedes the fitted $2.022$ of R2. We push the
assumption-free engine to $N=22$ ($10.2$~hours for that unit alone) and the
reduction to $N=32$ --- $4.29\times10^{9}$ basis states, $21$~GB, $815$~s, the
RAM ceiling. Where the two overlap the reduction is $10^{5}$ times faster and
gives the identical size multiset, and the closed form is exact at every $N$
either route reaches.
(iii) \emph{What is left open.} The sector is the maximal reachable set, not the
dynamics. $U$ has an exponentially large stationary space,
$\dim\Fix(U_w)=\binom{\lceil N/2\rceil}{w/2}$ and
$\dim\Fix(U)=2^{\lceil N/2\rceil}$ (an exact law, verified two independent
ways), almost all of it \emph{dark superpositions} that no basis state names. The
spectrum of $U_w$ is degenerate for \emph{every} shell with $|S_w|>N+1$ and
nondegenerate for the extremal ones, so outside those extremes
$\dim\mathcal{K}(x) < |S_w|$ without exception: an inner sector is never the
Krylov space of one of its basis states. Under monitoring the sector \emph{is}
filled uniformly (provably); without monitoring the diagonal ensemble occupies
between $6.6\%$ and $50\%$ of it. Finally, in the bond representation C150 is a
free-fermion Hadamard wall walk dressed by exactly two diagonal signs --- a
maximal nearest-neighbour wall interaction and a Jordan--Wigner string on the
coin --- and removing either one alone leaves the many-body spectrum
non-additive while removing both makes it additive to $4\times10^{-15}$. C150 is
therefore exactly charge-resolved but genuinely interacting inside each sector.
\end{abstract}

\tableofcontents

\section{Why this rule, and a naming note}

Throughout this series rules are named by Wolfram number with a \rt{W} prefix
(\rt{W150}); the present report follows the user-facing shorthand \emph{C150} for
the same object. The encoding is the standard one of Tier 1: the Wolfram number
is read as a $4$-tuple of local symbols indexed by the control pattern
$(x_{i-1},x_{i+1})$, with $\rt{I}$ the identity and $\rt{V}$ the Hadamard, so
\begin{equation}
  150 \;\longmapsto\; (\rt{I},\rt{V},\rt{V},\rt{I}),
  \qquad\text{i.e.}\qquad
  \text{Hadamard fires at site } i \iff x_{i-1}\oplus x_{i+1} = 1 .
  \label{eq:rule}
\end{equation}
The firing condition is exactly the update of the \emph{classical} additive
elementary rule $90$; C150 is the automaton that applies a Hadamard wherever ECA
$90$ would write a one. One cycle is the brick-wall product of a layer on the
even sites followed by a layer on the odd sites, controls always read from the
current bitmask. Boundaries: \rt{obc0} fixes $x_{-1}=x_N=0$, which is the case
treated here; the ring \rt{pbc} appears only for contrast in
\S\ref{sec:closed}.

C150 is singled out for three reasons. It is one of the two rules
($150$ and $105$) whose sector sizes R2 identified as binomials, so a closed form
was plausible; it is the rule whose \rt{obc0} sector sizes are used as a
regression fixture against the external HSF oracle in R1; and its fitted growth
exponent ($2.022$) sat awkwardly above the informational ceiling $2$, which
suggested that the fit was seeing a prefactor rather than a base. All three are
settled below.

\section{Reduction: the sector partition of a unitary rule is a flip graph}
\label{sec:flip}

\begin{theorem}[flip reduction]\label{thm:flip}
Let the rule be unitary (symbols in $\{\rt{I},\rt{V}\}$) and let
$s_i(x) = t[2x_{i-1}+x_{i+1}]$ be the symbol fired at site $i$ when the controls
read $x$. Then the sector partition --- the weakly connected components of the
one-cycle transition graph, which for a unitary rule coincide with its strongly
connected components --- equals the connected components of the undirected graph
\begin{equation}
  F:\quad x \;\sim\; x \oplus 2^i \qquad\text{for every } i
  \text{ with } s_i(x) = \rt{V}.
  \label{eq:flip}
\end{equation}
\end{theorem}

\begin{proof}
A cycle is two depth-$1$ half-layers. By the no-interference theorem of R1, each
half-layer is --- conditioned on the frozen complementary sublattice --- a
product of single-qubit gates acting on disjoint qubits, so the image of a basis
state $x$ under one half-layer is exactly the \emph{subcube} spanned by the sites
whose symbol is $\rt{V}$, with the complementary sublattice held fixed. Two
consequences. First, $x$ lies in its own half-layer image, because
$\bra{b}\rt{I}\ket{b}=1$ and $\bra{b}H\ket{b}=\pm1/\sqrt2$ are both nonzero;
hence the composite image $\Phi(x)$ contains both half-layer images $A(x)$ and
$B(x)$. Second, inside a half-layer image the frozen sublattice --- and therefore
the $\rt{V}/\rt{I}$ pattern itself --- does not change, so the subcube is
connected by single flips which are themselves edges of $F$. By the first point
every $F$-edge is realised by a composite path; by the second, every composite
edge $x\to z$ (with $z\in B(y)$, $y\in A(x)$) is realised by the $F$-path
$x\sim y \sim z$. The two graphs have the same components.
\end{proof}

Two remarks. The reduction is \emph{layer-agnostic}: the brick-wall ordering
disappears, because the controls of site $i$ are read off $x$ itself in the
half-layer in which $i$ is frozen. The sector partition of a unitary QECA is
therefore a property of the local symbol table alone --- not of the circuit
schedule. And the cost collapses from $O(2^N|\mathrm{succ}|)$, with
$|\mathrm{succ}|$ growing exponentially for the Hadamard-rich rules, to
$O(N2^N)$ with a tiny constant; \S\ref{sec:frontier} uses this.

\paragraph{Validation.} Theorem~\ref{thm:flip} is checked against the streamed
engine (\fn{graph.scc.analyze}, which knows nothing about it) for all $16$
unitary rules, both boundary conventions and $N=3\ldots12$: $320$ units, full
size multisets compared, no mismatch, for both the pure-python union-find and the
vectorised label-propagation implementation. See
\fn{tests/test\_flip\_graph.py}.

\section{Solution: walls, and the Ising charge}
\label{sec:walls}

Introduce the bond (domain-wall) variables. For \rt{obc0} there are $N+1$ bonds,
\begin{equation}
  b_j = x_j \oplus x_{j+1}, \quad j = 0,\dots,N,
  \qquad x_{-1} = x_N = 0 ,
\end{equation}
with the index convention of the code ($b_0 = x_0$, $b_N = x_{N-1}$).

\begin{proposition}[bond bijection]\label{prop:bij}
$x \mapsto b$ is a bijection from $\{0,1\}^N$ onto the even-weight strings of
length $N+1$; the inverse is the prefix XOR $x_j = \bigoplus_{k\le j} b_k$. On
the ring the same map is exactly $2$-to-$1$ onto the even-weight strings of
length $N$, the two preimages being $x$ and its global spin flip.
\end{proposition}
Both statements are verified exhaustively for $N = 2\ldots12$ (obc0) and
$N=3\ldots12$ (pbc).

Now translate the flip move \eqref{eq:flip}. Flipping $x_i$ toggles exactly the
two bonds $b_i$ and $b_{i+1}$, and by \eqref{eq:rule} the flip is allowed iff
$x_{i-1}\neq x_{i+1}$, i.e.\ iff $b_i \neq b_{i+1}$. So the elementary move is
\begin{equation}
  (b_i, b_{i+1}) = (1,0) \;\longleftrightarrow\; (0,1) :
  \qquad\text{a \emph{hard-core hop} of one domain wall.}
  \label{eq:hop}
\end{equation}
Creation and annihilation of walls is forbidden --- the configuration
$b_i = b_{i+1}$ is precisely where the symbol is $\rt{I}$. Therefore:

\begin{corollary}[conserved charge]\label{cor:charge}
The wall number $w=\sum_j b_j$ is conserved exactly, i.e.\ $[U,Q]=0$ for the
strictly $2$-local diagonal operator
\begin{equation}
  Q \;=\; \sum_{j=0}^{N} \frac{1 - Z_jZ_{j+1}}{2},
  \qquad Z_{-1} = Z_{N} = +1 ,
\end{equation}
the open Ising domain-wall number with boundary fields. ($Q$ is diagonal, so
conservation at the level of supports already gives the operator identity;
support-level conservation is verified for every $x$, $N=3\ldots12$, both
boundary conventions, with zero violations.)
\end{corollary}

\begin{corollary}[connectivity]\label{cor:conn}
Each wall shell is one sector. Hard-core hopping on a segment of $N+1$ sites is
connected on every fixed-occupancy shell (bubble the walls to the leftmost packed
configuration), and \eqref{eq:hop} is exactly that dynamics.
\end{corollary}

\section{Closed form, and the growth law}
\label{sec:closed}

Combining Proposition~\ref{prop:bij} with Corollaries~\ref{cor:charge}
and~\ref{cor:conn}:
\begin{equation}
  \boxed{\;
  |S_w| = \binom{N+1}{w}\ \ (w \text{ even}), \qquad
  \#\text{sectors} = \Big\lfloor \frac{N+1}{2}\Big\rfloor + 1, \qquad
  D_{\max} = \max_{w\ \mathrm{even}} \binom{N+1}{w}. \;}
  \label{eq:closed}
\end{equation}
The frozen sectors are the singletons $w=0$ (all-zero) and, when $N$ is odd,
$w=N+1$ (the alternating configuration); for even $N$ there is only one. The
parity bookkeeping in $D_{\max}$ is not cosmetic: for $N\equiv 1 \pmod 4$ the
central binomial sits at odd $w$ and is \emph{not} a sector size, so
$D_{\max}$ is the near-central one (e.g.\ $N=17$: $\binom{18}{8}=43758$, not
$\binom{18}{9}$).

\paragraph{Ring, for contrast.} On the ring the wall map is $2$-to-$1$, so a
shell with $0<w<N$ lifts to one sector of size $2\binom{N}{w}$, while the frozen
shells $w=0$ and (for even $N$) $w=N$ each split into \emph{two} singletons --- no
hop is available, and the two spin preimages are disconnected. Hence
$\#\text{sectors} = N/2+3$ for even $N$ and $(N+3)/2$ for odd $N$. Both closed
forms reproduce the stored Tier-1a data at every $N=6\ldots20$, in both boundary
conventions, size multiset for size multiset.

\paragraph{C150 is not fragmented.} Equation~\eqref{eq:closed} has a sector count
that is \emph{linear} in $N$ and sector sizes that are full eigenspaces of the
single charge $Q$. There is no Krylov structure beyond the $U(1)$ symmetry: C150
is \emph{symmetry-resolved}, not fragmented. This is worth stating plainly
because C150 is the rule most often quoted from this family, and it is exactly
the wrong example for fragmentation. The fragmented members of the family look
nothing like it --- \rt{W204} has $2^N$ frozen singletons, and the wall families
of \rt{W51}/\rt{W201} carry exponentially many sectors. C150 is the family's null
model.

\paragraph{The growth law: base exactly $2$, prefactor $N^{-1/2}$.}
Since $\sum_{w\,\mathrm{even}}\binom{N+1}{w} = 2^N$ and the maximum is (near-)
central,
\begin{equation}
  \frac{D_{\max}}{2^N} \;=\; 2\sqrt{\frac{2}{\pi (N+1)}}\,\bigl(1+O(N^{-1})\bigr)
  \;\longrightarrow\; 0 ,
  \label{eq:asym}
\end{equation}
so $D_{\max}$ grows as $2^N N^{-1/2}$: the exponential base is \emph{exactly} $2$
and the sub-exponential correction is a square root. The approach is not smooth
but a sawtooth (Figure~\ref{fig:frontier}, right), which is the parity
bookkeeping above: for $N\equiv1 \pmod 4$ the true central binomial is at odd $w$
and $D_{\max}$ has to settle for the neighbour, which costs a factor
$1-O(N^{-1})$ and produces the dips at $N=5,9,13,\ldots$. A pure-exponential fit
sees this sawtooth on top of the $N^{-1/2}$ decay, which is exactly the
configuration that pushes an unconstrained base slightly above the true one. The
value $2.022$ reported
for C150/\rt{obc0} in R2's pure-exponential fit is the finite-size shadow of that
prefactor, not a base above the informational ceiling; R2's companion observation
that $150/105$ carry $\alpha=-\tfrac12$ is the same statement seen from the fit
side, and \eqref{eq:asym} now supplies it exactly. Figure~\ref{fig:frontier}
(right) overlays the exact ratio on \eqref{eq:asym}.

\section{Numerical frontier}
\label{sec:frontier}

Two frontiers are reported, deliberately kept separate.

\begin{itemize}
\item The \textbf{assumption-free frontier}: the streamed engine
  (\fn{graph.scc.sectors\_union\_find} over \fn{core.cycle.succ}), which
  propagates exact $\mathbb{Z}[1/\sqrt2]$ amplitudes and uses no result of this
  report. Its cost grows by a factor $\approx 3$ per site because the one-cycle
  support itself grows exponentially, so it is the binding runtime constraint:
  $N=20$ already takes $65$ minutes.
\item The \textbf{reduction frontier}: Theorem~\ref{thm:flip} implemented as
  vectorised label propagation on the flip graph, with the label array
  ($\mathtt{uint32}$, $4$~bytes per basis state) the only object of size $2^N$ ---
  pointer jumping is done in place and in chunks, and the size histogram is
  harvested chunkwise, so nothing else of that size is ever allocated. This is
  RAM-bound rather than time-bound, and it reaches $N=32$: $2^{32} =
  4\,294\,967\,296$ basis states, $17$ sectors,
  $D_{\max} = 1\,166\,803\,110$, in $815$~s at a peak of $21.0$~GB. That is the
  ceiling twice over --- $N=32$ is the largest index a $\mathtt{uint32}$ label can
  hold, and $N=33$ would need a $64$-bit label array of $68.7$~GB.
\end{itemize}

Both agree with each other and with \eqref{eq:closed} at every $N$ computed.
Table~\ref{tab:frontier} lists the units; runtimes are wall-clock on the
project machine ($20$ cores, $31$~GB, jobs \fn{nice}d and run concurrently). The
two routes overlap at $N=21$ and $N=22$, where they return the identical size
multiset and the reduction is faster by $4.0\times10^{4}$ and
$1.0\times10^{5}$ respectively ($13\,018$~s versus $0.32$~s, $36\,728$~s versus
$0.37$~s). That ratio is the price of streaming the full one-cycle support
instead of a spanning set of edges, and it is what makes the difference between
$N=22$ and $N=32$.

\begin{table}[htbp]
\centering\footnotesize
\begin{tabular}{rrrrrll}
\hline
$N$ & $2^N$ & $\#$sectors & $D_{\max}$ & $D_{\max}/2^N$ & engine & flip graph \\
\hline
14 & $16\,384$ & $8$ & $6\,435$ & $0.3928$ & \cmark~4.4\,s & --- \\
15 & $32\,768$ & $9$ & $12\,870$ & $0.3928$ & \cmark~14\,s & --- \\
16 & $65\,536$ & $9$ & $24\,310$ & $0.3709$ & \cmark~40\,s & --- \\
17 & $131\,072$ & $10$ & $43\,758$ & $0.3338$ & \cmark~126\,s & --- \\
18 & $262\,144$ & $10$ & $92\,378$ & $0.3524$ & \cmark~363\,s & --- \\
19 & $524\,288$ & $11$ & $184\,756$ & $0.3524$ & \cmark~1321\,s & --- \\
20 & $1\,048\,576$ & $11$ & $352\,716$ & $0.3364$ & \cmark~1.1\,h & --- \\
21 & $2\,097\,152$ & $12$ & $646\,646$ & $0.3083$ & \cmark~3.6\,h & \cmark~0.3\,s \\
22 & $4\,194\,304$ & $12$ & $1\,352\,078$ & $0.3224$ & \cmark~10.2\,h & \cmark~0.4\,s \\
23 & $8\,388\,608$ & $13$ & $2\,704\,156$ & $0.3224$ & --- & \cmark~0.9\,s \\
24 & $16\,777\,216$ & $13$ & $5\,200\,300$ & $0.3100$ & --- & \cmark~1.8\,s \\
25 & $33\,554\,432$ & $14$ & $9\,657\,700$ & $0.2878$ & --- & \cmark~3.7\,s \\
26 & $67\,108\,864$ & $14$ & $20\,058\,300$ & $0.2989$ & --- & \cmark~8.9\,s \\
27 & $134\,217\,728$ & $15$ & $40\,116\,600$ & $0.2989$ & --- & \cmark~18\,s \\
28 & $268\,435\,456$ & $15$ & $77\,558\,760$ & $0.2889$ & --- & \cmark~38\,s \\
29 & $536\,870\,912$ & $16$ & $145\,422\,675$ & $0.2709$ & --- & \cmark~72\,s \\
30 & $1\,073\,741\,824$ & $16$ & $300\,540\,195$ & $0.2799$ & --- & \cmark~154\,s \\
31 & $2\,147\,483\,648$ & $17$ & $601\,080\,390$ & $0.2799$ & --- & \cmark~297\,s \\
32 & $4\,294\,967\,296$ & $17$ & $1\,166\,803\,110$ & $0.2717$ & --- & \cmark~815\,s \\
\hline
\end{tabular}

\caption{C150 \rt{obc0}: exact sector data and the two numerical frontiers.
\cmark{} means the computed size multiset equals the closed form
\eqref{eq:closed} entry for entry; the time is that unit's wall clock; ``---''
means that route has not been run at that $N$ (the engine is a factor
$\approx3$ per site, so its units above $N=20$ are hours to days each). The
closed-form columns are exact for every $N$, computed or not.}
\label{tab:frontier}
\end{table}

\begin{figure}[htbp]
\centering
\ifgraphic{../../figures/fig_c150_frontier.pdf}{\textwidth}
\caption{Left: the exact sector structure \eqref{eq:closed} with the two
computed frontiers marked. Right: $D_{\max}/2^N$ against the asymptotic
$2\sqrt{2/\pi(N+1)}$ of \eqref{eq:asym} --- the base is $2$ and the decay is the
$N^{-1/2}$ prefactor that a pure-exponential fit misreads as $2.022$.}
\label{fig:frontier}
\end{figure}

\section{What the sector does \emph{not} determine}
\label{sec:open}

A sector is the minimal $U$-invariant subspace spanned by basis states: the
\emph{maximal kinematically reachable} set from any of its members, and nothing
more (R1, \S``What a sector is''). For C150 we can now quantify the gap exactly,
because the shells are small enough to diagonalise. The dataset is every shell of
every $N=4\ldots24$ with $|S_w|\le3000$: $70$ shells, the largest of dimension
$2380$ ($N=16$, $w=4$). All numbers below come from $U_w$ built column by column
from the engine's exact amplitudes and verified orthogonal to $10^{-9}$. The
eigenvalue clustering tolerance is $10^{-9}$, and it has margin in both
directions: across the whole dataset the degeneracies of $U_w$ are exact
(largest gap treated as a degeneracy $8.8\times10^{-15}$, five orders below the
tolerance) while the smallest gap treated as \emph{distinct} is
$4.8\times10^{-6}$, more than three orders above it.

\paragraph{An exponentially large stationary space.} $U$ does not merely fail to
be ergodic inside a sector; it has a huge kernel of $U-\mathbb{1}$:
\begin{equation}
  \boxed{\;\dim\Fix(U_w) \;=\; \binom{\lceil N/2\rceil}{w/2},
  \qquad
  \dim\Fix(U) \;=\; 2^{\lceil N/2\rceil} \;=\; 2^{\#\{\text{even sites}\}} . \;}
  \label{eq:fix}
\end{equation}
The shell law is verified over the whole dataset --- every shell of every
$N=4\ldots24$ with $|S_w|\le3000$, $70$ shells, no mismatch; the total is
verified for $N=3\ldots12$ by two
independent routes, the multiplicity of the eigenvalue $+1$ and
$\dim\ker(U-\mathbb{1})$ by rank. The mechanism is visible in the half-layer
structure: writing $U = L_BL_A$ with $L_A,L_B$ the (Hermitian, involutive)
half-layers, one finds \emph{numerically for $N=3\ldots11$} that
\begin{equation}
  \Fix(U) \;=\; \Fix(L_A)\cap\Fix(L_B),
\end{equation}
i.e.\ a state stationary under the cycle is already stationary under each
half-layer separately, and each half-layer is a product of commuting single-qubit
involutions ($H$ or $\mathbb{1}$ per site of its sublattice), each of which
contributes a one-dimensional local $+1$ eigenspace. The counting exponent
$\lceil N/2\rceil$ is the number of even sites, i.e.\ the width of the first
half-layer. Only the $w=0$ (and, for odd $N$, $w=N+1$) members of $\Fix(U)$ are
basis states; \emph{all the rest are dark superpositions}, invisible to any
computational-basis graph.

\paragraph{Consequence: an inner sector is never a Krylov space.} For a basis
state $\ket{x}$, $\dim\mathcal{K}(x)$ equals the number of distinct eigenvalues of
$U_w$ that $\ket{x}$ overlaps, so any degeneracy bounds it strictly below
$|S_w|$. Over the whole dataset the dividing line is sharp and depends only on
the sector dimension:
\begin{equation}
\begin{aligned}
  \text{spectrum of } U_w \text{ nondegenerate}
  \;&\iff\; |S_w| \le N+1 \\
  \;&\iff\; w \in \{0,\,N+1\} \ \text{ or }\ (N \text{ even and } w = N),
\end{aligned}
\end{equation}
i.e.\ exactly the extremal shells --- the two frozen singletons and, for even
$N$, the top shell of dimension $N+1$, where \eqref{eq:fix} gives
$\dim\Fix(U_w)=1$ and no further multiplicity appears. For every one of the $59$
\emph{inner} shells ($|S_w|>N+1$) the spectrum is degenerate and
\begin{equation}
  \dim\mathcal{K}(x) \;<\; |S_w|,
  \qquad \frac{\dim\mathcal{K}(x)}{|S_w|} \in [0.395,\,0.963],
\end{equation}
the fraction decreasing with sector size (Table~\ref{tab:spectral},
Figure~\ref{fig:quantum}). This is the sharp form, for this rule, of the general
caveat in R1: the graph gives $\dim\mathcal{K}(x)\le|S_w|$, and away from the
extremal shells the inequality is always strict. For odd $N$ the reflection $P$
commutes with $U$ exactly ($[P,U]=0$ to machine zero, $30$ of the $70$ shells)
and reflection-symmetric seeds lose roughly half their Krylov space again
($N=7$, $w=2$: $13$ instead of $25$) --- the symmetry-forced degeneracy
anticipated in R1, here measured. For even $N$ the brick-wall order maps the
sublattices into each other and reflection acts instead as a time reversal,
$PUP = U^{-1}$ exactly (the other $40$ shells).

\paragraph{Monitored versus unmonitored, made numerical.} Under monitoring the
sector \emph{is} filled uniformly, provably: $M_w = |U_w|^2$ is doubly
stochastic, irreducible on $S_w$ (Corollary~\ref{cor:conn}) and aperiodic
(by the R1 no-interference theorem $\bra{x}U\ket{x}\neq0$ for every $x$, so every
node carries a self-loop), hence $M_w^t \to$ uniform on $S_w$. Without
monitoring the long-time state is the diagonal ensemble
$\bar\rho = \sum_\lambda P_\lambda\ket{x}\bra{x}P_\lambda$, and it is far from
uniform: over the $59$ inner shells its effective dimension
$d_{\mathrm{eff}} = 1/\operatorname{tr}\bar\rho^2$ occupies between $6.6\%$ and
$50\%$ of the sector --- never more than half --- and the most-occupied basis
state carries $2.0$ to $12.5$ times the uniform weight
(Table~\ref{tab:spectral}). So for C150 the two experiments genuinely disagree
about the long-time state \emph{inside} a sector, even though they agree exactly
about which sector that state lives in, and (by R1) about populations after a
single cycle.

\begin{table}[htbp]
\centering\footnotesize\setlength{\tabcolsep}{4pt}
\begin{tabular}{rrrrrrrrrl}
\hline
$N$ & $w$ & $|S_w|$ & $\#\mathrm{spec}$ & $\dim\mathrm{Fix}$ & $\binom{\lceil N/2\rceil}{w/2}$ & $\dim\mathcal{K}/|S_w|$ & $d_{\mathrm{eff}}/|S_w|$ & $\max p\cdot|S_w|$ & reflection \\
\hline
8 & 2 & $36$ & $33$ & $4$ & $4$ & $0.917$ & $0.388$ & $2.57$ & $PUP=U^{-1}$ \\
8 & 4 & $126$ & $81$ & $6$ & $6$ & $0.643$ & $0.222$ & $4.46$ & $PUP=U^{-1}$ \\
8 & 6 & $84$ & $65$ & $4$ & $4$ & $0.774$ & $0.292$ & $3.42$ & $PUP=U^{-1}$ \\
10 & 2 & $55$ & $51$ & $5$ & $5$ & $0.927$ & $0.359$ & $2.78$ & $PUP=U^{-1}$ \\
10 & 4 & $330$ & $211$ & $10$ & $10$ & $0.639$ & $0.156$ & $6.33$ & $PUP=U^{-1}$ \\
10 & 6 & $462$ & $243$ & $10$ & $10$ & $0.526$ & $0.137$ & $7.17$ & $PUP=U^{-1}$ \\
10 & 8 & $165$ & $131$ & $5$ & $5$ & $0.794$ & $0.241$ & $4.12$ & $PUP=U^{-1}$ \\
10 & 10$^\dagger$ & $11$ & $11$ & $1$ & $1$ & $1.000$ & $0.793$ & $1.26$ & $PUP=U^{-1}$ \\
12 & 2 & $78$ & $73$ & $6$ & $6$ & $0.936$ & $0.337$ & $2.96$ & $PUP=U^{-1}$ \\
12 & 4 & $715$ & $473$ & $15$ & $15$ & $0.662$ & $0.123$ & $8.07$ & $PUP=U^{-1}$ \\
12 & 6 & $1\,716$ & $729$ & $20$ & $20$ & $0.425$ & $0.079$ & $12.49$ & $PUP=U^{-1}$ \\
12 & 8 & $1\,287$ & $665$ & $15$ & $15$ & $0.517$ & $0.096$ & $10.15$ & $PUP=U^{-1}$ \\
12 & 10 & $286$ & $233$ & $6$ & $6$ & $0.815$ & $0.218$ & $4.58$ & $PUP=U^{-1}$ \\
12 & 12$^\dagger$ & $13$ & $13$ & $1$ & $1$ & $1.000$ & $0.785$ & $1.27$ & $PUP=U^{-1}$ \\
14 & 2 & $105$ & $99$ & $7$ & $7$ & $0.943$ & $0.316$ & $3.13$ & $PUP=U^{-1}$ \\
14 & 4 & $1\,365$ & $939$ & $21$ & $21$ & $0.688$ & $0.106$ & $9.37$ & $PUP=U^{-1}$ \\
14 & 12 & $455$ & $379$ & $7$ & $7$ & $0.833$ & $0.201$ & $4.98$ & $PUP=U^{-1}$ \\
14 & 14$^\dagger$ & $15$ & $15$ & $1$ & $1$ & $1.000$ & $0.779$ & $1.28$ & $PUP=U^{-1}$ \\
16 & 2 & $136$ & $129$ & $8$ & $8$ & $0.949$ & $0.312$ & $3.20$ & $PUP=U^{-1}$ \\
16 & 4 & $2\,380$ & $1\,697$ & $28$ & $28$ & $0.713$ & $0.093$ & $10.75$ & $PUP=U^{-1}$ \\
16 & 14 & $680$ & $577$ & $8$ & $8$ & $0.849$ & $0.191$ & $5.08$ & $PUP=U^{-1}$ \\
16 & 16$^\dagger$ & $17$ & $17$ & $1$ & $1$ & $1.000$ & $0.774$ & $1.29$ & $PUP=U^{-1}$ \\
\hline
\end{tabular}

\caption{Wall shells of C150 \rt{obc0}. $\#\mathrm{spec}$ is the number of
distinct eigenvalues of $U_w$; $\dim\Fix$ its $+1$ multiplicity, beside the
closed form \eqref{eq:fix}; $\dim\mathcal{K}/|S_w|$ the largest Krylov fraction
over the sampled seeds; $d_{\mathrm{eff}}/|S_w|$ the smallest diagonal-ensemble
fraction; $\max p\cdot|S_w|$ the peak long-time population in units of the
uniform value ($1$ would be uniform, which is what monitoring gives exactly).}
\label{tab:spectral}
\end{table}

\begin{figure}[htbp]
\centering
\ifgraphic{../../figures/fig_c150_quantum.pdf}{\textwidth}
\caption{Left: the Krylov and diagonal-ensemble fractions of a sector against
sector dimension; the dashed line is the monitored answer, which is exactly
uniform. Filled markers are inner shells; open markers the extremal shells
$|S_w|\le N+1$ --- the only ones
with a nondegenerate spectrum, and hence the only ones a basis state's Krylov
space fills. Right: the stationary space
$\dim\Fix(U) = 2^{\lceil N/2\rceil}$, whose members other than the frozen basis
states are dark superpositions.}
\label{fig:quantum}
\end{figure}

\section{C150 is charge-resolved but interacting}
\label{sec:free}

The wall picture invites the guess that C150 is a free-fermion circuit. It is
not, and the obstruction is exactly two signs.

Write the cycle on the $N+1$ bond modes. The gate at site $i$ touches bonds $i$
and $i+1$ only; it is the identity when both carry the same occupation, and on
the one-wall subspace --- basis ordered (wall at bond $i$, wall at bond $i+1$) ---
it is
\begin{equation}
  M(c) \;=\; \frac{1}{\sqrt2}
  \begin{pmatrix} -(-1)^c & 1 \\ 1 & (-1)^c \end{pmatrix},
  \qquad \det M(c) = -1 ,
  \label{eq:coin}
\end{equation}
where $c = x_{i-1} = \bigoplus_{k<i} b_k$ is the parity of the walls strictly to
the left of bond $i$ --- a Jordan--Wigner string. This bond-space circuit
reproduces the engine's $U$ under the bijection of Proposition~\ref{prop:bij} to
$1.1\times10^{-16}$ (Table~\ref{tab:freeness}, second column), so
\eqref{eq:coin} \emph{is} C150: a hard-core quantum walk of domain walls with a
Hadamard coin.

Two features of \eqref{eq:coin} spoil gaussianity. A number-conserving Gaussian
two-mode gate with single-particle block $u$ must act on the doubly occupied
state by $\det u$; here $\det M = -1$ while the gate acts by $+1$, which is a
maximal nearest-neighbour interaction $\exp(i\pi n_i n_{i+1})$ between walls. And
the $(-1)^c$ string multiplies the \emph{density-difference} part of the coin,
not the hop, so it is not absorbed by the Jordan--Wigner transform. The test is
sharp: a number-conserving Gaussian circuit acts on the $w$-particle sector as
the $w$-th antisymmetric power of its single-particle block, so its eigenphases
are subset sums of the one-particle phases. Table~\ref{tab:freeness} reports the
largest failure of that additivity over $w=2\ldots5$ for the four sign variants.

\begin{table}[htbp]
\centering\footnotesize
\begin{tabular}{rrrrrr}
\hline
$N$ & bond vs engine & rule 150 & no string & $\det$ phase only & Gaussian ref. \\
\hline
5 & $1.1\times10^{-16}$ & $1.871$ & $1.101$ & $1.340$ & $8.9\times10^{-16}$ \\
6 & $1.1\times10^{-16}$ & $1.157$ & $0.869$ & $0.869$ & $1.4\times10^{-15}$ \\
7 & $1.1\times10^{-16}$ & $1.062$ & $0.608$ & $0.526$ & $1.4\times10^{-15}$ \\
8 & $1.1\times10^{-16}$ & $0.731$ & $0.688$ & $0.515$ & $3.0\times10^{-15}$ \\
9 & $1.1\times10^{-16}$ & $0.549$ & $0.604$ & $0.427$ & $3.7\times10^{-15}$ \\
\hline
\end{tabular}

\caption{Additivity defect: the largest $|\arg(\lambda/\lambda')|$ in radians
between the $w$-wall spectrum and the $w$-fold products of the one-wall
eigenvalues, maximised over $w=2\ldots5$ (compared as points on the unit circle,
so the branch cut cannot masquerade as a defect). Column 2 is the bond
representation checked against the engine. C150 (column 3) and each of the two
single-defect models (columns 4, 5) are non-additive at $0.43$--$1.87$~rad;
removing \emph{both} defects gives the Gaussian reference (column 6), additive to
$4\times10^{-15}$.}
\label{tab:freeness}
\end{table}

The conclusion is a clean separation of the two questions. The
\emph{fragmentation} of C150 is entirely a $U(1)$ symmetry and is solved in
closed form; the \emph{dynamics within a sector} is that of interacting hard-core
walls, and the interaction is not a perturbation --- it is maximal. That is the
honest reason a sector's spectral properties do not follow from
\eqref{eq:closed}, and why \eqref{eq:fix} had to be measured rather than read off
the graph.

\section{Open items}

\begin{enumerate}
\item \textbf{Level statistics.} Deliberately not computed here. The spectra of
  every shell in Table~\ref{tab:spectral} are reproducible in seconds from
  \fn{quantum.rule150\_spectra}, so a level-spacing analysis --- sector by
  sector, and separated by reflection parity for odd $N$ --- is immediately
  available; \S\ref{sec:free} predicts it will \emph{not} look free, and
  \eqref{eq:fix} says the $+1$ eigenvalue must be excised first.
\item \textbf{A derivation of \eqref{eq:fix}.} The law
  $\dim\Fix(U_w)=\binom{\lceil N/2\rceil}{w/2}$ and the identity
  $\Fix(U)=\Fix(L_A)\cap\Fix(L_B)$ are exact numerical facts here, with the
  half-layer mechanism identified but not turned into a proof. The natural route
  is the local $+1$ eigenvector of each single-gate involution, together with the
  count of even sites; the $w$-resolved binomial suggests the dark states are
  labelled by choosing $w/2$ of the $\lceil N/2\rceil$ even sites.
\item \textbf{The further exact multiplicities.} Beyond the $+1$ eigenvalue the
  spectrum carries exact multiplicities $2,3,4,6,10,15$ (e.g.\ $N=8$, $w=4$:
  $\{1{:}48,\,2{:}24,\,3{:}8,\,6{:}1\}$). These are not explained by reflection,
  which is abelian; some further commutant is present.
\item \textbf{The ring.} \S\ref{sec:closed} solves \rt{pbc} kinematically, but
  the quantum layer of \S\ref{sec:open}--\S\ref{sec:free} is done for \rt{obc0}
  only. On the ring the wall map is $2$-to-$1$, so the bond representation
  acquires a $\mathbb{Z}_2$ gauge redundancy and \eqref{eq:coin} needs a flux
  sector; that is the next natural step.
\item \textbf{Higher $N$ at the assumption-free frontier.} The streamed engine
  costs a factor $\approx2.8$ per site ($N=20$: $1.1$~h; $N=21$: $3.6$~h;
  $N=22$: $10.2$~h), so $N=23$ is a $\approx30$-hour unit --- launched, but
  outside the one-day budget of this report, and not needed for the closed form.
  The reduction frontier is instead
  bounded by the label array: $N=33$ needs $\mathtt{uint64}$ and $68.7$~GB, so
  the next step there is a machine with more memory, or a compact relabelling
  (the components are few, so a two-pass scheme with a byte-wide label after the
  first sweep would halve the requirement).
\end{enumerate}

\section{Reproduce}

\begin{verbatim}
# closed forms, wall bijection, charge conservation, engine cross-check
python -m qca_fragmentation.c150 --verify

# assumption-free frontier (one unit; N=22 takes ~10 h)
python -m qca_fragmentation.run_rule --rule 150 --N 21-22 --bc obc0

# reduction frontier (checkpointed to analytics/c150_frontier.jsonl)
python -m qca_fragmentation.c150 --frontier 21-31 --bc obc0

# sector spectra, Krylov dimensions, Fix(U), freeness test
python -m qca_fragmentation.scaling.c150_report --quantum
python -m qca_fragmentation.scaling.c150_report --tables --figures

# tests (126 of them)
cd code/qca_fragmentation/tests
python -m pytest test_flip_graph.py test_c150.py -q
\end{verbatim}

\begin{sloppypar}
Data: \fn{results/150\_obc0.jsonl} and \fn{results/150\_pbc.jsonl} (Tier 1a),
\fn{analytics/c150\_frontier.jsonl} (reduction frontier, one JSON line per
$(N,\rt{bc})$ with the full size multiset and the closed-form check),
\fn{analytics/c150\_quantum.json} (spectral portraits, $\Fix(U)$, freeness).
\end{sloppypar}

\end{document}
